\documentclass[a4paper, 12pt]{article}

\usepackage{utils}

\renewcommand*{\today}{20 January 2026}

\begin{document}

\hotbox{Questions ouvertes}{CM 1}{\today}

Programme fait en python pour tester en partant d'un chiffre avec un multiple donné
\begin{lstlisting}
def until_turned(first_value, mul, limit=100):
    value = [first_value]
    result = []
    retenu = 0
    for i in range(limit):
        if result and retenu == 0 and result[-1] == first_value:
            break
        num = (value[-1] * mul + retenu) % 10
        retenu_tmp = (value[-1] * mul + retenu) // 10
        result.append(num)
        value.append(num)
        retenu = retenu_tmp
    return (''.join(map(str, value[::-1][1:])), len(value) - 1, ''.join(map(str, result[::-1])), len(result))

mul = 6
for i in range(1, 10):
    print(f"pour {i} avec mul = {mul} on a {until_turned(i, mul)}")
\end{lstlisting}

\begin{remarque}
    Le nombre de chiffres dépend du chiffre multipliant de départ.
    \begin{itemize}
        \item pour 2 on a 18 chiffres
        \item pour 3 on a 28 chiffres
        \item pour 4 on a 6 chiffres
        \item pour 5 on a 42 chiffres sauf 7 x 5
        \item pour 6 on a 58 chiffres
        \item pour 7 on a 22 chiffres
        \item pour 8 on a 13 chiffres
        \item pour 9 on a 44 chiffres
        \item pour 10 on a 2 chiffres
    \end{itemize}
\end{remarque}



\end{document}